% ===== Appendix: corrections to the book, what we added, sources (after Questions; not presented) =====
\appendix
\partaccent{lavender}\gdef\idiottopicemoji{books}\gdef\idiottopicname{Հավելված}
\newcolumntype{L}[1]{>{\raggedright\arraybackslash}p{#1}}

% Corrections in deck order, split over two frames since 2026-11-12 (six rows did not fit one)
\begin{frame}{Հավելված Ա. Ուղղումներ գրքի նկատմամբ (1/2)}
  \scriptsize % not presented; at \footnotesize the longer rows ran off the bottom
  \setlength{\tabcolsep}{4pt}
  \renewcommand{\arraystretch}{1.25}
  \begin{tabular}{@{}L{2.4cm}L{3.5cm}L{4.45cm}L{3.2cm}@{}}
    \toprule
    \textbf{Թեմա} & \textbf{Գրքում (2016)} & \textbf{Այսօր} & \textbf{Աղբյուր}\\
    \midrule
    «Flashbulb» հիշողություններ & Տրավմայի վառ, կրկնվող հիշողություններ & Տերմինը ցնցող լուրն իմանալու պահի վառ հիշողությունն է, որը ճշգրիտ չէ. տրավմայի կրկնվողները՝ ինտրուզիվ հիշողություններ (PTSD) & Brown \& Kulik, 1977; Talarico \& Rubin, 2003\\
    Սուզորդների փորձը & Նույն միջավայրը լավացնում է վերհիշումը & 16 սուզորդով կրկնությունն էֆեկտը չգտավ & Murre, R. Soc. Open Sci., 2021\\
    Բլեքաութ (blackout) & Լրիվ բլեքաութը մեծ մասի մոտ լինում է, երբ հազիվ է կանգնում. նորմալ խոսելով՝ միայն կախվածություն ունեցողների մոտ & Խմող ուսանողների կեսը գոնե մեկ անգամ բլեքաութ է ունեցել. դրսից համարյա չի երևում & White et al., J. Am. Coll. Health, 2002; NIAAA\\
    \bottomrule
  \end{tabular}
\end{frame}

\begin{frame}{Հավելված Ա. Ուղղումներ գրքի նկատմամբ (2/2)}
  \scriptsize % not presented; at \footnotesize the longer rows ran off the bottom
  \setlength{\tabcolsep}{4pt}
  \renewcommand{\arraystretch}{1.25}
  \begin{tabular}{@{}L{2.4cm}L{3.5cm}L{4.45cm}L{3.2cm}@{}}
    \toprule
    \textbf{Թեմա} & \textbf{Գրքում (2016)} & \textbf{Այսօր} & \textbf{Աղբյուր}\\
    \midrule
    Ջոն Դին & Ձայնագրությունները ցույց տվեցին, որ նա «լավագույն դեպքում երկրորդական» դեր ուներ & Դինը քողարկման կենտրոնում էր ու մեղավորությունն ընդունեց. այդ զրույցներում իր դերը հիշում էր ավելի կենտրոնական, քան էր & Neisser, Cognition, 1981\\
    Պացիենտ Հ.Մ. & Մահացել է 78 տարեկանում & 82 տարեկանում (1926-2008) & Squire, Neuron, 2009\\
    Ալցհեյմեր & Գլխավոր տեսությունը՝ նեյրոֆիբրիլների խճճվածքները & Ամիլոիդի սալիկներ (դրսում) ու tau-ի խճճվածքներ (ներսում). 2023-ից ամիլոիդը մաքրող դեղերը մի քիչ դանդաղեցնում են հիվանդությունը & van Dyck et al., NEJM, 2023; Sims et al., JAMA, 2023\\
    \bottomrule
  \end{tabular}
\end{frame}

\begin{frame}{Հավելված Բ. Ինչ ենք ավելացրել գրքից դուրս}
  \begin{columns}[T]
    \begin{column}{0.48\textwidth}
      \begin{card}[colback=mint!30!cardwhite, fontupper=\scriptsize]
        {\blackfont Թվեր ու հետազոտություններ}\par\smallskip
        \begin{itemize}\setlength{\itemsep}{0.5mm}
          \item «Քուն» բառի թեստը (DRM)՝ Roediger \& McDermott, 1995
          \item Լոֆթուսի թվերը՝ Loftus \& Palmer, 1974
          \item Տաքսու վարորդների հիպոկամպը (հետին մասը)՝ Maguire et al., 2000
          \item Հ.Մ.-ի հիպոկամպերի մոտ կեսը՝ Corkin et al., J. Neurosci., 1997
          \item Ուսանողների բլեքաութները՝ White et al., 2002; NIAAA
          \item Ալցհեյմերի դեղերը՝ van Dyck et al., 2023; Sims et al., 2023
          \item 140\%-ն ու զույգերը՝ Veritasium, 2020
        \end{itemize}
      \end{card}
    \end{column}
    \begin{column}{0.48\textwidth}
      \begin{card}[colback=peach!35!cardwhite, fontupper=\scriptsize]
        {\blackfont Մերն է, ոչ թե գրքից}\par\smallskip
        \begin{itemize}\setlength{\itemsep}{0.5mm}
          \item Անուն հիշելու 2-րդ ու 3-րդ խորհուրդները
          \item Chunking-ի օրինակները (թվերն ու կատուն). գրքում՝ մեկ այլ նախադասություն
          \item Հ.Մ.-ի հայելու փորձի մանրամասները (Բրենդա Միլներ). գրքում՝ «նկարչական հմտություններ»
          \item Մեմը. բովանդակությունը գրքից է
          \item Խմբային նախագծի հարցումը, «Առասպել, թե՞ փաստ»-ը, անցումային հարցերը
        \end{itemize}
      \end{card}
    \end{column}
  \end{columns}
\end{frame}

\begin{frame}{Հավելված Գ. Աղբյուրներ. գլուխ 2}
  \scriptsize
  Գրքի ծանոթագրություններից, որոնց վրա հիմնված են սլայդները\par\smallskip
  \begin{columns}[T]
    \begin{column}{0.48\textwidth}
      \begin{itemize}\setlength{\itemsep}{0.4mm}
        \item Cowan N. The magical mystery four. Curr. Dir. Psychol. Sci., 2010
        \item Miller G. A. The magical number seven, plus or minus two. Psychol. Rev., 1956 (գիրքը հղում է Nicolis \& Tsuda, 1985-ին)
        \item Godden D. R., Baddeley A. D. Context-dependent memory in two natural environments. Br. J. Psychol., 1975
        \item Henson R. N. Short-term memory for serial order. Cogn. Psychol., 1998
        \item Eichenbaum H. The Cognitive Neuroscience of Memory. 2011
        \item Greenwald A. G. The totalitarian ego. Am. Psychol., 1980
        \item Neisser U. John Dean's memory: a case study. Cognition, 1981
      \end{itemize}
    \end{column}
    \begin{column}{0.48\textwidth}
      \begin{itemize}\setlength{\itemsep}{0.4mm}
        \item Mather M., Johnson M. K. Choice-supportive source monitoring. Psychol. Aging, 2000
        \item Meissner C. A., Brigham J. C. Thirty years of investigating the own-race bias. Psychol. Public Policy Law, 2001
        \item Hoffrage U. et al. Hindsight bias: a by-product of knowledge updating? JEP:LMC, 2000
        \item Walker W. R., Skowronski J. J. The fading affect bias. Appl. Cogn. Psychol., 2009
        \item Wilson A. E., Ross M. From chump to champ. JPSP, 2001
        \item Loftus E. F. The price of bad memories. 1998 (Նադին Քուլի դեպքը)
        \item Squire L. R. The legacy of patient H.M. for neuroscience. Neuron, 2009
      \end{itemize}
    \end{column}
  \end{columns}
\end{frame}

% Alcohol and the brain (Hayk asked, 2026-11-12): beyond the book, checked against the sources named on each card
\newcommand{\alccard}[4]{% {colour}{title}{text}{source}
  \begin{card}[colback=#1!35!cardwhite, fontupper=\footnotesize]
    {\blackfont\small #2}\par\smallskip
    #3\par\smallskip
    {\color{plum!65}#4}
  \end{card}}

% four pages of two cards (2026-11-12): eight cards on two pages ran 70-95 pt off the bottom
\begin{frame}{Հավելված Դ. Ալկոհոլն ու ուղեղը (1/4)}
  \begin{columns}[T]
    \begin{column}{0.48\textwidth}
      \alccard{mint}{Արգելակն ու գազը}{Ալկոհոլը ուժեղացնում է GABA-ն՝ ուղեղի գլխավոր «արգելակը», ու թուլացնում գլուտամատը՝ գլխավոր «գազը»։ Հիպոկամպում գլուտամատի NMDA ընկալիչները պետք են նոր հիշողություն գրանցելու համար, ու ալկոհոլն արգելափակում է հենց դրանք։}{Lovinger et al., Science, 1989; White, Alcohol Res. Health, 2003}
    \end{column}
    \begin{column}{0.48\textwidth}
      \alccard{butter}{Քունը}{Ալկոհոլով ավելի արագ ես քնում, ու գիշերվա առաջին կեսն ավելի խորն է։ Բայց REM քունն ուշանում է, իսկ երկրորդ կեսն ավելի անհանգիստ է։ Իսկ քնի ժամանակ է ուղեղն ամրացնում օրվա հիշողությունները։}{Ebrahim et al., Alcohol Clin. Exp. Res., 2013}
    \end{column}
  \end{columns}
\end{frame}

\begin{frame}{Հավելված Դ. Ալկոհոլն ու ուղեղը (2/4)}
  \begin{columns}[T]
    \begin{column}{0.48\textwidth}
      \alccard{lavender}{Խմել սովորելուց հետո}{88 մարդ սովորել է բառեր, հետո կեսը խմել է, կեսը՝ ոչ։ Հաջորդ առավոտ խմածներն ավելի լավ էին հիշում, ու ինչքան շատ էին խմել, այնքան լավ։ Հավանական պատճառը՝ ալկոհոլը փակում է նոր տեղեկության մուտքը, ու հինն ամրանում է առանց «խանգարողների»։ Մեկ փոքր հետազոտություն է, ու քունն էլ կարող էր ազդել։}{Carlyle et al., Sci. Rep., 2017}
    \end{column}
    \begin{column}{0.48\textwidth}
      \alccard{peach}{Կարծում ես՝ խմել ես, ուրեմն գեղեցիկ ես}{Տղամարդիկ, որոնք կարծում էին, թե ալկոհոլ են խմել, իրենց ավելի գրավիչ էին գնահատում, նույնիսկ եթե իրականում պլացեբո էր։ Անկախ գնահատողները տարբերություն չտեսան։ Ազդում էր ոչ թե դոզան, այլ սպասելիքը։}{Bègue et al., Br. J. Psychol., 2013}
    \end{column}
  \end{columns}
\end{frame}

\begin{frame}{Հավելված Դ. Ալկոհոլն ու ուղեղը (3/4)}
  \begin{columns}[T]
    \begin{column}{0.48\textwidth}
      \alccard{mint}{Նեյրոնները չի «սպանում», բայց...}{Ալկոհոլը, հավանաբար, ուղղակիորեն չի սպանում նեյրոնները։ Բայց առնետների մոտ երկու շաբաթ արյան մոտ 0.08\% ալկոհոլով (ԱՄՆ-ում մեքենա վարելու սահմանը) հիպոկամպում նոր նեյրոնների ստեղծումը մոտ 40\%-ով պակասեց։}{Anderson et al., Neuroscience, 2012; The Conversation}
    \end{column}
    \begin{column}{0.48\textwidth}
      \alccard{lavender}{Ուղեղի ծավալը}{36\,678 չափահասի ուղեղի MRI (Մեծ Բրիտանիա). ծավալն ավելի փոքր էր արդեն օրական 1-2 միավորի դեպքում (1 միավոր = 8 գ ալկոհոլ, մոտ 250 մլ գարեջուր)։ Հեղինակների հաշվով՝ օրական 2-ից 3 միավորի անցնելը նման է 3.5 տարով ծերանալուն։ Կապ է, ոչ ապացուցված պատճառ։}{Daviet et al., Nat. Commun., 2022}
    \end{column}
  \end{columns}
\end{frame}

\begin{frame}{Հավելված Դ. Ալկոհոլն ու ուղեղը (4/4)}
  \begin{columns}[T]
    \begin{column}{0.48\textwidth}
      \alccard{peach}{Թեթև խմելն էլ չի պաշտպանում}{Դիտարկումներում թվում էր, թե քիչ խմողների մոտ դեմենցիան ավելի հազվադեպ է։ 2025-ին 2.4 մլն մարդու գենետիկ տվյալները ցույց տվեցին՝ ռիսկն աճում է խմածի հետ միասին։ «Պաշտպանությունը», հավանաբար, հակառակ կապ էր. դեմենցիան սկսվելիս մարդիկ ավելի քիչ են խմում։}{Topiwala et al., BMJ Evid. Based Med., 2025}
    \end{column}
    \begin{column}{0.48\textwidth}
      \alccard{butter}{Ինչու է վտանգավոր հանկարծ թողնելը}{Երկար ու շատ խմելիս ուղեղը հարմարվում է. ուժեղացնում է «գազը» (գլուտամատ), որ ալկոհոլին հակակշռի։ Հանկարծ թողնելիս հակակշիռը մնում է՝ դող, անհանգստություն, ծանր դեպքում ցնցումներ։ Դրա համար ծանր կախվածությունից դուրս են գալիս բժշկի հսկողությամբ։}{Littleton, Alcohol Health Res. World, 1998; Becker, 1998}
    \end{column}
  \end{columns}
\end{frame}
